% Options for packages loaded elsewhere
% Options for packages loaded elsewhere
\PassOptionsToPackage{unicode}{hyperref}
\PassOptionsToPackage{hyphens}{url}
\PassOptionsToPackage{dvipsnames,svgnames,x11names}{xcolor}
%
\documentclass[
  letterpaper,
  DIV=11,
  numbers=noendperiod]{scrartcl}
\usepackage{xcolor}
\usepackage{amsmath,amssymb}
\setcounter{secnumdepth}{-\maxdimen} % remove section numbering
\usepackage{iftex}
\ifPDFTeX
  \usepackage[T1]{fontenc}
  \usepackage[utf8]{inputenc}
  \usepackage{textcomp} % provide euro and other symbols
\else % if luatex or xetex
  \usepackage{unicode-math} % this also loads fontspec
  \defaultfontfeatures{Scale=MatchLowercase}
  \defaultfontfeatures[\rmfamily]{Ligatures=TeX,Scale=1}
\fi
\usepackage{lmodern}
\ifPDFTeX\else
  % xetex/luatex font selection
\fi
% Use upquote if available, for straight quotes in verbatim environments
\IfFileExists{upquote.sty}{\usepackage{upquote}}{}
\IfFileExists{microtype.sty}{% use microtype if available
  \usepackage[]{microtype}
  \UseMicrotypeSet[protrusion]{basicmath} % disable protrusion for tt fonts
}{}
\makeatletter
\@ifundefined{KOMAClassName}{% if non-KOMA class
  \IfFileExists{parskip.sty}{%
    \usepackage{parskip}
  }{% else
    \setlength{\parindent}{0pt}
    \setlength{\parskip}{6pt plus 2pt minus 1pt}}
}{% if KOMA class
  \KOMAoptions{parskip=half}}
\makeatother
% Make \paragraph and \subparagraph free-standing
\makeatletter
\ifx\paragraph\undefined\else
  \let\oldparagraph\paragraph
  \renewcommand{\paragraph}{
    \@ifstar
      \xxxParagraphStar
      \xxxParagraphNoStar
  }
  \newcommand{\xxxParagraphStar}[1]{\oldparagraph*{#1}\mbox{}}
  \newcommand{\xxxParagraphNoStar}[1]{\oldparagraph{#1}\mbox{}}
\fi
\ifx\subparagraph\undefined\else
  \let\oldsubparagraph\subparagraph
  \renewcommand{\subparagraph}{
    \@ifstar
      \xxxSubParagraphStar
      \xxxSubParagraphNoStar
  }
  \newcommand{\xxxSubParagraphStar}[1]{\oldsubparagraph*{#1}\mbox{}}
  \newcommand{\xxxSubParagraphNoStar}[1]{\oldsubparagraph{#1}\mbox{}}
\fi
\makeatother

\usepackage{color}
\usepackage{fancyvrb}
\newcommand{\VerbBar}{|}
\newcommand{\VERB}{\Verb[commandchars=\\\{\}]}
\DefineVerbatimEnvironment{Highlighting}{Verbatim}{commandchars=\\\{\}}
% Add ',fontsize=\small' for more characters per line
\usepackage{framed}
\definecolor{shadecolor}{RGB}{241,243,245}
\newenvironment{Shaded}{\begin{snugshade}}{\end{snugshade}}
\newcommand{\AlertTok}[1]{\textcolor[rgb]{0.68,0.00,0.00}{#1}}
\newcommand{\AnnotationTok}[1]{\textcolor[rgb]{0.37,0.37,0.37}{#1}}
\newcommand{\AttributeTok}[1]{\textcolor[rgb]{0.40,0.45,0.13}{#1}}
\newcommand{\BaseNTok}[1]{\textcolor[rgb]{0.68,0.00,0.00}{#1}}
\newcommand{\BuiltInTok}[1]{\textcolor[rgb]{0.00,0.23,0.31}{#1}}
\newcommand{\CharTok}[1]{\textcolor[rgb]{0.13,0.47,0.30}{#1}}
\newcommand{\CommentTok}[1]{\textcolor[rgb]{0.37,0.37,0.37}{#1}}
\newcommand{\CommentVarTok}[1]{\textcolor[rgb]{0.37,0.37,0.37}{\textit{#1}}}
\newcommand{\ConstantTok}[1]{\textcolor[rgb]{0.56,0.35,0.01}{#1}}
\newcommand{\ControlFlowTok}[1]{\textcolor[rgb]{0.00,0.23,0.31}{\textbf{#1}}}
\newcommand{\DataTypeTok}[1]{\textcolor[rgb]{0.68,0.00,0.00}{#1}}
\newcommand{\DecValTok}[1]{\textcolor[rgb]{0.68,0.00,0.00}{#1}}
\newcommand{\DocumentationTok}[1]{\textcolor[rgb]{0.37,0.37,0.37}{\textit{#1}}}
\newcommand{\ErrorTok}[1]{\textcolor[rgb]{0.68,0.00,0.00}{#1}}
\newcommand{\ExtensionTok}[1]{\textcolor[rgb]{0.00,0.23,0.31}{#1}}
\newcommand{\FloatTok}[1]{\textcolor[rgb]{0.68,0.00,0.00}{#1}}
\newcommand{\FunctionTok}[1]{\textcolor[rgb]{0.28,0.35,0.67}{#1}}
\newcommand{\ImportTok}[1]{\textcolor[rgb]{0.00,0.46,0.62}{#1}}
\newcommand{\InformationTok}[1]{\textcolor[rgb]{0.37,0.37,0.37}{#1}}
\newcommand{\KeywordTok}[1]{\textcolor[rgb]{0.00,0.23,0.31}{\textbf{#1}}}
\newcommand{\NormalTok}[1]{\textcolor[rgb]{0.00,0.23,0.31}{#1}}
\newcommand{\OperatorTok}[1]{\textcolor[rgb]{0.37,0.37,0.37}{#1}}
\newcommand{\OtherTok}[1]{\textcolor[rgb]{0.00,0.23,0.31}{#1}}
\newcommand{\PreprocessorTok}[1]{\textcolor[rgb]{0.68,0.00,0.00}{#1}}
\newcommand{\RegionMarkerTok}[1]{\textcolor[rgb]{0.00,0.23,0.31}{#1}}
\newcommand{\SpecialCharTok}[1]{\textcolor[rgb]{0.37,0.37,0.37}{#1}}
\newcommand{\SpecialStringTok}[1]{\textcolor[rgb]{0.13,0.47,0.30}{#1}}
\newcommand{\StringTok}[1]{\textcolor[rgb]{0.13,0.47,0.30}{#1}}
\newcommand{\VariableTok}[1]{\textcolor[rgb]{0.07,0.07,0.07}{#1}}
\newcommand{\VerbatimStringTok}[1]{\textcolor[rgb]{0.13,0.47,0.30}{#1}}
\newcommand{\WarningTok}[1]{\textcolor[rgb]{0.37,0.37,0.37}{\textit{#1}}}

\usepackage{longtable,booktabs,array}
\newcounter{none} % for unnumbered tables
\usepackage{calc} % for calculating minipage widths
% Correct order of tables after \paragraph or \subparagraph
\usepackage{etoolbox}
\makeatletter
\patchcmd\longtable{\par}{\if@noskipsec\mbox{}\fi\par}{}{}
\makeatother
% Allow footnotes in longtable head/foot
\IfFileExists{footnotehyper.sty}{\usepackage{footnotehyper}}{\usepackage{footnote}}
\makesavenoteenv{longtable}
\usepackage{graphicx}
\makeatletter
\newsavebox\pandoc@box
\newcommand*\pandocbounded[1]{% scales image to fit in text height/width
  \sbox\pandoc@box{#1}%
  \Gscale@div\@tempa{\textheight}{\dimexpr\ht\pandoc@box+\dp\pandoc@box\relax}%
  \Gscale@div\@tempb{\linewidth}{\wd\pandoc@box}%
  \ifdim\@tempb\p@<\@tempa\p@\let\@tempa\@tempb\fi% select the smaller of both
  \ifdim\@tempa\p@<\p@\scalebox{\@tempa}{\usebox\pandoc@box}%
  \else\usebox{\pandoc@box}%
  \fi%
}
% Set default figure placement to htbp
\def\fps@figure{htbp}
\makeatother





\setlength{\emergencystretch}{3em} % prevent overfull lines

\providecommand{\tightlist}{%
  \setlength{\itemsep}{0pt}\setlength{\parskip}{0pt}}



 


\KOMAoption{captions}{tableheading}
\makeatletter
\@ifpackageloaded{caption}{}{\usepackage{caption}}
\AtBeginDocument{%
\ifdefined\contentsname
  \renewcommand*\contentsname{Table of contents}
\else
  \newcommand\contentsname{Table of contents}
\fi
\ifdefined\listfigurename
  \renewcommand*\listfigurename{List of Figures}
\else
  \newcommand\listfigurename{List of Figures}
\fi
\ifdefined\listtablename
  \renewcommand*\listtablename{List of Tables}
\else
  \newcommand\listtablename{List of Tables}
\fi
\ifdefined\figurename
  \renewcommand*\figurename{Figure}
\else
  \newcommand\figurename{Figure}
\fi
\ifdefined\tablename
  \renewcommand*\tablename{Table}
\else
  \newcommand\tablename{Table}
\fi
}
\@ifpackageloaded{float}{}{\usepackage{float}}
\floatstyle{ruled}
\@ifundefined{c@chapter}{\newfloat{codelisting}{h}{lop}}{\newfloat{codelisting}{h}{lop}[chapter]}
\floatname{codelisting}{Listing}
\newcommand*\listoflistings{\listof{codelisting}{List of Listings}}
\makeatother
\makeatletter
\makeatother
\makeatletter
\@ifpackageloaded{caption}{}{\usepackage{caption}}
\@ifpackageloaded{subcaption}{}{\usepackage{subcaption}}
\makeatother
\usepackage{bookmark}
\IfFileExists{xurl.sty}{\usepackage{xurl}}{} % add URL line breaks if available
\urlstyle{same}
\hypersetup{
  pdftitle={Lab: Power Calculations},
  colorlinks=true,
  linkcolor={blue},
  filecolor={Maroon},
  citecolor={Blue},
  urlcolor={Blue},
  pdfcreator={LaTeX via pandoc}}


\title{Lab: Power Calculations}
\usepackage{etoolbox}
\makeatletter
\providecommand{\subtitle}[1]{% add subtitle to \maketitle
  \apptocmd{\@title}{\par {\large #1 \par}}{}{}
}
\makeatother
\subtitle{Group Project Protocol Guidelines}
\author{}
\date{}
\begin{document}
\maketitle


\subsection{Introduction}\label{introduction}

The goal of this exercise is to attempt a power calculation based for an
experiment. This is based on the chosen design and linear model for the
experiment, for instance a Complete Block Design, Crossover Design etc.
We shall first run through an example using the Complete Block Design
setup. You will then need to run a similar power calculation for your
group's chosen design and setup and submit it as part of the written
Protocol.

\subsection{Complete Block Design
Example}\label{complete-block-design-example}

We now look at an example calculation for a sample Complete Blocked
Design \(CB[1]\) experiment in \texttt{R}. Note that we use fixed Block
effects here as opposed to random effects.

\paragraph{Step 1: Linear Model and Data Generating
Process}\label{step-1-linear-model-and-data-generating-process}

We first identify the linear model associated with our design as
follows:

\[
Y_i = \mu + \alpha_{j(i)} + \beta_{k(i)} + \epsilon_i
\]

\[
\hat{\mu} = \bar{Y}
\]

\[
\hat{\alpha}_{j(i)} = \bar{Y}_{j(i)} - \hat{\mu}
\]

\[
\hat{\beta}_{k(i)} = \bar{Y}_{k(i)} - \hat{\mu}
\]

\textbf{Where:}

\begin{itemize}
\tightlist
\item
  \(Y_i\) is the response for observation \(i\);\\
\item
  \(\mu\) is the overall (grand) mean;
\item
  \(\alpha_{j(i)}\) is the factor effects, where \(j(i)\) indicates the
  factor level of the \(i\)th observation;
\item
  \(\beta_{k(i)}\) are block effects, where \(k(i)\) indicates the block
  the \(i\)th observation belongs to;\\
\item
  \(\epsilon_i\) is the error term,
  \(\epsilon_i \overset{\text{i.i.d.}}{\sim} \mathcal{N}(0, \sigma^2)\).
\end{itemize}

\(\bar{Y}, \bar{Y}_{j(i)}, \bar{Y}_{k(i)}\) are sample means.\\

Write this out in R as follows:

\begin{Shaded}
\begin{Highlighting}[]
\CommentTok{\# Data Generating Process outline}

\CommentTok{\# Write out the linear model as an R function}

\NormalTok{generate\_Y }\OtherTok{\textless{}{-}} \ControlFlowTok{function}\NormalTok{(mu, alpha, beta, sigma\_sq, j, k) \{}
  \CommentTok{\# j: n length vector of factor level indices}
  \CommentTok{\# k: n length vector of block indices}
  
\NormalTok{  n }\OtherTok{\textless{}{-}} \FunctionTok{length}\NormalTok{(j)}
  
  \FunctionTok{set.seed}\NormalTok{(}\DecValTok{42}\NormalTok{)}
\NormalTok{  epsilon }\OtherTok{\textless{}{-}} \FunctionTok{rnorm}\NormalTok{(n, }\AttributeTok{mean =} \DecValTok{0}\NormalTok{, }\AttributeTok{sd =} \FunctionTok{sqrt}\NormalTok{(sigma\_sq))}
  
\NormalTok{  Y }\OtherTok{\textless{}{-}}\NormalTok{ mu }\SpecialCharTok{+}\NormalTok{ alpha[j] }\SpecialCharTok{+}\NormalTok{ beta[k] }\SpecialCharTok{+}\NormalTok{ epsilon}
  
  \FunctionTok{return}\NormalTok{(Y)}
\NormalTok{\}}
\end{Highlighting}
\end{Shaded}

\paragraph{Step 2: Choosing
Parameters}\label{step-2-choosing-parameters}

We now choose the relevant parameters for our power calculation. Set
effect sizes to the smallest value that would be scientifically
meaningful to you. For instance, . You may also choose these from
previous literature etc. Keep in mind that your size should match
conditions on these parameters (sum to zero etc)!

Next, use an estimate of the error variance \(\sigma^2\), ideally from
your pilot data or dry run. Set the sample size as a meaningful sample
size that you feasibly hope to attain or collect while conducting your
experiment.

For the purposes of our example we stick to two factor levels and two
blocks, i.e 4 observations in total.

\begin{Shaded}
\begin{Highlighting}[]
\CommentTok{\# Defining parameter values}

\NormalTok{mu }\OtherTok{\textless{}{-}}
\NormalTok{alpha }\OtherTok{\textless{}{-}} \FunctionTok{c}\NormalTok{(,) }\CommentTok{\# Fill in levels of alpha: they should sum to zero!}
\NormalTok{beta }\OtherTok{\textless{}{-}} \FunctionTok{c}\NormalTok{(,)  }\CommentTok{\# Fill in levels of beta: they should sum to zero!}

\NormalTok{sigma\_sq }\OtherTok{\textless{}{-}}
\NormalTok{n }\OtherTok{\textless{}{-}} 
\end{Highlighting}
\end{Shaded}

Note in particular the role of \(\sigma^2\) in the power calculations.
Once you finish steps 1-4, remember to come back and double the value of
\(\sigma^2\). What do you see?

\paragraph{Step 3: Simulating a dataset and
Inference}\label{step-3-simulating-a-dataset-and-inference}

We now simulate a dataset using the linear data-generating process and
the parameters defined above. We can do this simply as follows.

Fill in the missing steps, in particular choose \(j\) and \(k\) for each
observation and store it in the vectors. Sketching out a table for your
experiment may be useful here!

\begin{Shaded}
\begin{Highlighting}[]
\CommentTok{\# Simulating the dataset}

\NormalTok{j }\OtherTok{\textless{}{-}} \FunctionTok{c}\NormalTok{()  }\CommentTok{\# fill in factor levels}
\NormalTok{k }\OtherTok{\textless{}{-}} \FunctionTok{c}\NormalTok{()  }\CommentTok{\# fill in blocks}

\CommentTok{\# Simulate response using the model}
\NormalTok{Y }\OtherTok{\textless{}{-}} \FunctionTok{generate\_Y}\NormalTok{(mu, alpha, beta, sigma\_sq, j, k)}

\CommentTok{\# Put into a data frame}
\NormalTok{dat }\OtherTok{\textless{}{-}} \FunctionTok{data.frame}\NormalTok{(}
  \AttributeTok{Y =}\NormalTok{ Y,}
  \AttributeTok{treatment =} \FunctionTok{factor}\NormalTok{(j),}
  \AttributeTok{block =} \FunctionTok{factor}\NormalTok{(k)}
\NormalTok{)}
\end{Highlighting}
\end{Shaded}

We now proceed to run inference and calculate the p-value for our single
factor here using the traditional ANOVA and F-test. Complete the steps!

\begin{Shaded}
\begin{Highlighting}[]
\CommentTok{\# Inference on simulated data }

\CommentTok{\# Fit linear model (fixed effects for both)}
\NormalTok{fit }\OtherTok{\textless{}{-}} \FunctionTok{lm}\NormalTok{(Y }\SpecialCharTok{\textasciitilde{}} \CommentTok{\# fill in linear model}
\NormalTok{          , }\AttributeTok{data =}\NormalTok{ dat)}

\CommentTok{\# Extract ANOVA table}
\NormalTok{aov\_tab }\OtherTok{\textless{}{-}} \FunctionTok{anova}\NormalTok{(fit)}

\CommentTok{\# Get p{-}value for treatment effect}
\NormalTok{p\_val }\OtherTok{\textless{}{-}}\NormalTok{ aov\_tab[}\StringTok{"treatment"}\NormalTok{, }\StringTok{"Pr(\textgreater{}F)"}\NormalTok{]}

\CommentTok{\# Convert to rejection indicator}
\NormalTok{reject }\OtherTok{\textless{}{-}} \FunctionTok{as.integer}\NormalTok{(p\_val }\SpecialCharTok{\textless{}} \FloatTok{0.05}\NormalTok{)}

\NormalTok{reject}
\end{Highlighting}
\end{Shaded}

\paragraph{Step 4: Iterating to calculate
Power}\label{step-4-iterating-to-calculate-power}

Now that we have outlined the pipeline for simulating a dataset and
doing inference, we compile it together into a single \texttt{R}
function so that it is easy to replicate.

\begin{Shaded}
\begin{Highlighting}[]
\CommentTok{\# One simulation run through of Steps 1{-}3}

\NormalTok{one\_sim }\OtherTok{\textless{}{-}} \ControlFlowTok{function}\NormalTok{(mu, alpha, beta, sigma\_sq, j, k) \{}
  \CommentTok{\# Generate data}
\NormalTok{  Y }\OtherTok{\textless{}{-}} \FunctionTok{generate\_Y}\NormalTok{(mu, alpha, beta, sigma\_sq, j, k)}
  
\NormalTok{  dat }\OtherTok{\textless{}{-}} \FunctionTok{data.frame}\NormalTok{(}
    \AttributeTok{Y =}\NormalTok{ Y,}
    \AttributeTok{treatment =} \FunctionTok{factor}\NormalTok{(j),}
    \AttributeTok{block =} \FunctionTok{factor}\NormalTok{(k)}
\NormalTok{  )}
  
  \CommentTok{\# Fit model}
\NormalTok{  fit }\OtherTok{\textless{}{-}} \FunctionTok{lm}\NormalTok{(Y }\SpecialCharTok{\textasciitilde{}}\NormalTok{ treatment }\SpecialCharTok{+}\NormalTok{ block, }\AttributeTok{data =}\NormalTok{ dat)}
  
  \CommentTok{\# Extract ANOVA p{-}value for treatment}
\NormalTok{  p\_val }\OtherTok{\textless{}{-}} \FunctionTok{anova}\NormalTok{(fit)[}\StringTok{"treatment"}\NormalTok{, }\StringTok{"Pr(\textgreater{}F)"}\NormalTok{]}
  
  \CommentTok{\# Return rejection indicator}
  \FunctionTok{return}\NormalTok{(}\FunctionTok{as.integer}\NormalTok{(p\_val }\SpecialCharTok{\textless{}} \FloatTok{0.05}\NormalTok{))}
\NormalTok{\}}
\end{Highlighting}
\end{Shaded}

Now, we use our single-run function to iterate over many such simulated
datasets under the same proposed ``Alternate Hypothesis'' of our chosen
effect sizes. We use the Monte Carlo machinery from previous power
calculations to estimate the probability of rejecting the null () under
this alternate!

\begin{Shaded}
\begin{Highlighting}[]
\CommentTok{\# Iterate over Steps 1{-}3 multiple times: Turn the Monte Carlo crank!}

\CommentTok{\# Number of simulations}
\NormalTok{iter }\OtherTok{\textless{}{-}} \DecValTok{1000}

\CommentTok{\# Run simulations}
\NormalTok{reject\_vec }\OtherTok{\textless{}{-}} \FunctionTok{replicate}\NormalTok{(}
\NormalTok{  iter,}
  \FunctionTok{one\_sim}\NormalTok{(mu, alpha, beta, sigma\_sq, j, k)}
\NormalTok{)}

\CommentTok{\# Estimated power}
\NormalTok{power\_est }\OtherTok{\textless{}{-}} \FunctionTok{mean}\NormalTok{(reject\_vec)}

\NormalTok{power\_est}
\end{Highlighting}
\end{Shaded}

Thus, \texttt{power\_est} gives us an estimate for power under our
specified design and alternative (the minimum deviation from the null
that we care about!). We can also run some diagnostics to make sure this
estimate has correctly converged.

\begin{Shaded}
\begin{Highlighting}[]
\FunctionTok{library}\NormalTok{(ggplot2)}

\CommentTok{\# Running Monte Carlo estimate of power}
\NormalTok{running\_power }\OtherTok{\textless{}{-}} \FunctionTok{cumsum}\NormalTok{(reject\_vec) }\SpecialCharTok{/} \FunctionTok{seq\_along}\NormalTok{(reject\_vec)}

\CommentTok{\# Put into a data frame for ggplot}
\NormalTok{df\_power }\OtherTok{\textless{}{-}} \FunctionTok{data.frame}\NormalTok{(}
  \AttributeTok{iter =} \FunctionTok{seq\_along}\NormalTok{(reject\_vec),}
  \AttributeTok{power =}\NormalTok{ running\_power}
\NormalTok{)}

\CommentTok{\# Convergence plot}
\FunctionTok{ggplot}\NormalTok{(df\_power, }\FunctionTok{aes}\NormalTok{(}\AttributeTok{x =}\NormalTok{ iter, }\AttributeTok{y =}\NormalTok{ power)) }\SpecialCharTok{+}
  \FunctionTok{geom\_line}\NormalTok{(}\AttributeTok{linewidth =} \FloatTok{0.8}\NormalTok{) }\SpecialCharTok{+}
  \FunctionTok{geom\_hline}\NormalTok{(}\AttributeTok{yintercept =} \FunctionTok{mean}\NormalTok{(reject\_vec), }\AttributeTok{linetype =} \StringTok{"dashed"}\NormalTok{) }\SpecialCharTok{+}
  \FunctionTok{labs}\NormalTok{(}
    \AttributeTok{title =} \StringTok{"Monte Carlo Power Estimate Convergence"}\NormalTok{,}
    \AttributeTok{x =} \StringTok{"Number of Simulations"}\NormalTok{,}
    \AttributeTok{y =} \StringTok{"Estimated Power"}
\NormalTok{  ) }\SpecialCharTok{+}
  \FunctionTok{theme\_minimal}\NormalTok{()}
\end{Highlighting}
\end{Shaded}

If the value of the estimate starts to stay constant after a few
iterations, it means that it is stable and reliable as an estimate of
your power!

Feel free to go back and change your parameters and rerun the power
analysis. In particular, investigate the effect of changing
\(\sigma^2\)!

\hfill\break
\#\# Applying this to your own Group Projects

Use the same pipeline as above, following Steps 1-4 for your own group
projects!

\textbf{Step 1:} Update the linear model and the data generating process
according to the design you choose for your experiment. Be careful about
random effects, and model constraints.

\textbf{Step 2:} Choose the effect sizes, and sample size that best
reflects your setting and the question you want to answer. Keep in mind
zero-sum constraints etc!

\textbf{Step 3:} Fill out the design matrix or table for your
experiment, and use this to fill out \texttt{i} and \texttt{j} vectors.
Add other factors or terms as necessary. Simulate a dataset and
hypothesis test.

\textbf{Step 4:} Iterate through these steps as many times as needed to
ensure your power estimate has converged!\\

This should give you an idea of what power you can achieve with your
chosen design and sample size. Discuss as a group if it is sufficient
for your purposes or not. A power of 0.7 - 0.9 is usually considered
ideal.

Experiment with different sample sizes to see what size is suitable for
your experiment. Feel free to try out different designs, random vs fixed
block effects, or play around with the number of factors, levels etc.
See which of these gives you the highest power, you may use this to pin
down a design/factors if you're still uncertain. Another way of
increasing power is by decreasing \(\sigma^2\), either by blocking,
protocols, or better measurement.

Evaluate your group project ideas and submit a protocol. Best of luck
with your experiments :)




\end{document}
